% Source: https://artofproblemsolving.com/community/c3h2384465_underlined_numbers
% Stable identity: aops-topic-2384465-q01
Let $P$ be the product of all ten numbers.

\textbf{At most two.} If $y$ is underlined, then $y$ equals the product of the other nine numbers. Multiplying both sides by $y$ gives $y^2=P$. The equation $y^2=P$ has at most two real solutions ($\pm\sqrt P$), and the ten numbers are different, so at most $2$ numbers are underlined.

\textbf{Two is possible.} Take
\[7,\ -7,\ 2,\ \dfrac12,\ 3,\ \dfrac13,\ 4,\ \dfrac14,\ 5,\ -\dfrac15.\]
These are ten different numbers, and the last eight have product $1\cdot1\cdot1\cdot(-1)=-1$, so $P=7\cdot(-7)\cdot(-1)=49$. The product of the other nine numbers is $\dfrac{P}{y}$ for each nonzero $y$: for $7$ it is $\dfrac{49}{7}=7$ and for $-7$ it is $-7$, so both are underlined, while no other number satisfies $y^2=49$.

So the answer is $\boxed{2}$.
